\begin{abstract}
We report a camera-driven underwater obstacle-avoidance study in which a
simulation calibrated against measurements of the target vehicle still
failed to transfer to it, and we trace why. Three calibration levels were
fitted to paired observations of a BlueROV2 and a test pool -- a static
observation model, a timing and burst-dropout model, and a vehicle
response model applied downstream of the planner -- and stacked on a
HoloOcean baseline. Eighty closed-loop runs covering two planners, two
obstacle geometries and five repetitions per cell were executed and
hashed \emph{before} any physical run, so that the field campaign would
test a prediction rather than illustrate one. The stack then did not
work on the vehicle. Making it work required seven documented changes,
after which nine runs were flown and eight are admissible, all with the
committed planner; the dynamic-window baseline never issued a command
because it needs a pose estimate the vehicle has no sensor for. Two
findings dominate, and neither is a calibrated quantity. The detector's
blob-size gates admit an obstacle only between \SI{0.29}{\metre} and
\SI{2.11}{\metre}, while the vehicle started \SI{1.86}{\metre} from the
anchor and engaged at \SI{1.80}{\metre}: the physical experiment had no
approach phase at all, against two metres of run-up in simulation. And
the planner's manoeuvre asks for a \SI{2.5}{\metre} lateral offset in a
basin \SIrange{2.7}{3.0}{\metre} wide. Because no external tracking was
recorded, the physical runs can validate the command domain and nothing
spatial, which we state as a result about what camera-only field
campaigns can prove rather than as a caveat. We recover the real
perception stream offline from the detector's own annotated video, audit
our own frozen campaign and report the defects we found in it, and use
post-deployment simulations -- explicitly labelled as such throughout --
to separate the effect of each deployment change, which the single
available pool day could not.
\keywords{Sim-to-real transfer \and Underwater robotics \and Obstacle
avoidance \and Monocular perception \and Deployment gap \and HoloOcean.}
\end{abstract}
